% =====================================================================
%  1 UŽDUOTIS — IoT tinklams būdingos kibernetinės atakos
%  Rugpjūčio 3–4 d. | Tikslinė apimtis: 4–5 psl.
% =====================================================================

\subsection{IoT architektūros sluoksniai ir jų atakų paviršius}

% TODO (rugpj. 3): trumpas trijų sluoksnių pristatymas (2 pastraipos)

\subsection{Atakų klasifikacija}

% TODO (rugpj. 3): po 1–2 pastraipas kiekvienam sluoksniui.
% Svarbiausia — stulpelis „matomi tinklo požymiai": jis tiesiogiai
% virsta modelio įvesties požymiais 4 skyriuje (\cref{sec:sprendimas}).

\begin{table}[H]
\centering
\caption{IoT atakų klasifikacija pagal architektūros sluoksnius}
\label{tab:atakos}
\small
\begin{tabularx}{\textwidth}{@{}l l X@{}}
\toprule
\textbf{Sluoksnis} & \textbf{Ataka} & \textbf{Matomi tinklo požymiai} \\
\midrule
Suvokimo & Jutiklių duomenų klastojimas & Netipinės telemetrijos reikšmės, staigūs šuoliai \\
         & Fizinis įsilaužimas & Naujas MAC adresas, pasikeitęs įrenginio profilis \\
\addlinespace
Tinklo   & DDoS / DoS & Didelis paketų dažnis, vienodas paketų dydis, daug SYN \\
         & MITM / ARP spoofing & Pasikartojantys ARP atsakymai, pasikeitęs maršrutas \\
         & Sybil & Daug tapatybių iš vieno fizinio šaltinio \\
         & RPL maršrutizavimo atakos & Anomalūs DIO/DAO pranešimai, topologijos kaita \\
\addlinespace
Programų & Mirai tipo botnetai & Telnet/SSH skenavimas, C\&C periodinės jungtys \\
         & MQTT piktnaudžiavimas & Netipiniai topic'ai, per didelis publish dažnis \\
         & Duomenų nutekinimas & Neįprastai didelis išeinantis srautas \\
\bottomrule
\end{tabularx}
\end{table}

\subsection{Klasikiniai aptikimo metodai ir jų ribos IoT kontekste}

% TODO (rugpj. 4): parašais grįstas (Snort/Suricata), taisyklėmis grįstas,
% anomalijomis grįstas — po pastraipą. Akcentas: KODĖL nepakanka.

\begin{table}[H]
\centering
\caption{Aptikimo metodų palyginimas IoT kontekste}
\label{tab:aptikimas}
\small
\begin{tabularx}{\textwidth}{@{}l X X@{}}
\toprule
\textbf{Metodas} & \textbf{Privalumai} & \textbf{Trūkumai IoT aplinkoje} \\
\midrule
Parašais grįstas & Tikslus žinomoms atakoms, mažai klaidingų teigiamų
& Neaptinka zero-day; parašų bazė reikalauja nuolatinio atnaujinimo \\
\addlinespace
Taisyklėmis / slenksčiais & Paprasta realizuoti, skaidru
& Daug klaidingų pavojaus signalų; slenksčiai netinka nevienalyčiam IoT elgesiui \\
\addlinespace
Anomalijomis grįstas & Aptinka nematytas atakas
& Reikia „normalaus elgesio" etalono; jautrus koncepcijos dreifui \\
\bottomrule
\end{tabularx}
\end{table}

\subsection{IoT specifika, ribojanti tradicinių IDS taikymą}

% TODO (rugpj. 4): riboti resursai, protokolų įvairovė, įrenginių mastas,
% nevienalytis normalaus elgesio profilis, šifruotas srautas.
% Ši dalis pagrindžia, kodėl 2 skyriuje kreipiamės į DI metodus.

\subsection{Aptikimo diegimo vieta}

% TODO (rugpj. 4): on-device / edge-gateway / cloud — privalumai ir trūkumai.
% Susieti su inferencijos delsos ir modelio dydžio matavimais 5 skyriuje.

\subsection{Skyriaus apibendrinimas}

% TODO: 1 pastraipa — kokie reikalavimai keliami aptikimo sprendimui.
